\section{Ablaufsplan und Projektorganisation}
\subsection{Vorgehen}

% Notizen. Quellen: Protokolle in docs/protokolle, internes Discord-Archiv, Workflow-Leitfaden im Repo.
% Literatur nur fuer methodische Konzepte; Projektdaten bleiben ohne Beleg.

\begin{itemize}
    \item Vorgehensmodell: iterativ-inkrementell in vier Phasen (Verstehen, Fundament, MVP, Integration) statt Wasserfall \autocite{larman2003iterative,sommerville2016software}
    \item Anforderungen empirisch erheben statt annehmen: Patientenbefragung, Therapeutenbefragung, Experteninterview, Literatur- und App-Vergleich \autocite{pohl2010requirements}
    \item Triangulationsregel: eine Anforderung gilt erst als gesichert, wenn mindestens zwei der Quellen sie stützen \autocite{flick2011triangulation}
    \item Priorisierung der Anforderungen in MUSS, KANN und OPTIONAL; alle sechs MUSS-Anforderungen sind umgesetzt
    \item frühe Prototypen als Diskussionsgrundlage: Wireframes in Figma, JSX-Prototyp, Netlify-Vorschau
    \item drei Teilgruppen entlang der Architektur: Knowledge Graph, API und UI/UX; ab 11.06. auch als Discord-Rollen sichtbar
    \item Entwicklungsprozess nach eigenem Workflow-Leitfaden im Repository:
    \begin{itemize}
        \item kein direkter Push auf den main-Branch, main bleibt immer lauffähig
        \item ein Branch pro Issue, Zusammenführung über Pull Requests \autocite{chacon2014progit}
        \item Review durch mindestens eine zweite Person vor dem Merge
        \item Labels für Gruppen (Graph, UI, API, Alle) und Typen (Feature, Bug, Design, Docs)
        \item Kanban-Board mit den Spalten Todo, In Progress, Review und Closed \autocite{anderson2010kanban}
    \end{itemize}
    \item Kommunikation: wöchentliches Freitagstreffen, ab Juni hybrid (Präsenz plus Discord); Protokolle im Repository; zusätzliche Gruppentreffen; Entscheidungen per Umfrage und Abstimmung
    \item Werkzeuge: GitHub, Discord, Nextcloud, Figma (Wireframes, Logo), Flutter/Dart, Isar, Protégé, Fuseki/SPARQL, Netlify (Vorschau), Canva
    \item Qualitätssicherung und gelernte Lektionen:
    \begin{itemize}
        \item Related Works am 22.05. beanstandet (Quellen nicht auffindbar), danach manuelle Überarbeitung
        \item Git-Integration Ende Juni bis Anfang Juli: UI entwickelte auf eigenem Stand, Merge-Konflikte, danach Einigung auf main als gemeinsame Basis
        \item Isar läuft nicht im Web, deshalb getrenntes Netlify-Vorschauprojekt
        \item ab 07.09. tägliche Kurzupdates der UI-Gruppe und wöchentliche Statusabfragen
    \end{itemize}
\end{itemize}

\subsection{Arbeitspakete}

\begin{itemize}
    \item AP1 Projektmanagement und Infrastruktur: Repository (GitLab ab 08.05., Umzug zu GitHub am 18./19.05.), Workflow-Leitfaden, Issues und Board, Protokolle (Sophia), Discord-Rollen (Hiep), Repo-Administration (Louis)
    \item AP2 Related Work und Anforderungen: Arbeitsgruppen A (Fragebogen), B (Fokus Mensch), C (Fokus Anwendung); Ergebnis: Anforderungskatalog mit MUSS, KANN und OPTIONAL
    \item AP3 Nutzerstudie: Fragebogenentwurf und Überarbeitung, Befragung in der BG Ambulanz (11.05. bis 29.05.), Fragebögen für vier Therapeut:innen, Experteninterview mit Etta Wolters (21.05.), Auswertung (Elif, Ilker)
    \item AP4 UI/UX-Design: Wireframes in Figma (Sophia, Nadine), Logo (Nadine), Design-System, Körperdiagramm
    \item AP5 MVP-Feature-Konzept: Ableitung aus Interviews, Fragebögen und Related Work, vorgestellt am 22.05. (Ilker)
    \item AP6 Wissensbasis: Ontologie in Protégé auf Basis von YogaKG und PhysicalTherapy.owx, Übungskatalog, SPARQL-Endpunkt (Knowledge-Gruppe)
    \item AP7 App und API: Flutter-Projekt ab 05.06., Nutzerdaten und Persistenz mit Isar, SPARQL-Anbindung, Planerstellung und Planempfehlung (API- und UI-Gruppe)
    \item AP8 Integration und Test: erster E2E-Durchstich am 22.06., Merge-Fixes im Juli, Betrieb auf dem Testgerät Galaxy S10 Lite ab 04.09., Medienanbindung im September
    \item AP9 Bericht und Präsentation: Gliederung nach Vorschlag der Betreuerin (02.07.), Beiträge pro Seite in der Fußzeile, Zwischenpräsentation am 26.06., MVP-Meilenstein am 17.07., Abschlusspräsentation am 02.10.
\end{itemize}

\subsection{Zeitlicher Ablauf}

% Soll: Projektplan vom 25.04. (docs/protokolle/REHAB-Uebersicht-Sophia.pdf).
% Planungsdokumente der Betreuerin (24./25.04.) sind im Discord-Archiv abgelaufen.

\textbf{Soll (Projektplan vom 25.04.):}
\begin{itemize}
    \item 17.04. Kick-off und Gruppeneinteilung, ab 24.04. Planung der Nutzerstudie
    \item 30.04. Fragebogen und Related Work, 08.05. Präsentation der Gruppenergebnisse
    \item 11.05. bis 22.05. Durchführung der Nutzerstudie in den Praxen
    \item Juni und Juli: MVP bauen (App-Gerüst, Backend-API, erste Queries, Demo mit kleinem Übungsset)
    \item August und September: Ausbau und Bewertung (Wissensbasis, Fortschritt, UI/UX, Erinnerungen)
    \item Abschluss: Bericht, Präsentationen, lauffähige Demo und Demo-Video
\end{itemize}

\textbf{Ist (vier Phasen):}
\begin{itemize}
    \item Phase 1 Verstehen (17.04. bis 26.05.): Kick-off, Gruppen, Related Work; Fragebögen raus am 11.05., Verlängerung der Befragung um eine Woche (18.05.); Interview am 21.05.; MVP-Feature-Ableitung am 22.05.; Meilenstein: Anforderungen belegt
    \item Phase 2 Fundament (04.06. bis 26.06.): Flutter-Start, Ontologie in Protégé, erste Queries, E2E-Test am 22.06.; Meilenstein: Zwischenpräsentation am 26.06.
    \item Phase 3 MVP (Ende Juni bis 19.07.): Git-Merge-Krise (30.06. bis 10.07.), Isar-Persistenz, MVP-Meilensteinpräsentation am 17.07., Feedback am 19.07.
    \item Phase 4 Integration (31.07. bis 02.10.): Planung der letzten Schritte am 31.07., Medien bis 15.09., Testgerät ab 04.09., Bericht bis 30.09., Abschlusspräsentation am 02.10.
\end{itemize}

\textbf{Abweichungen vom Plan:}
\begin{itemize}
    \item Nutzerstudie verlängert (Krankheit in der Praxis, größere Stichprobe)
    \item Related Works nach der Kritik am 22.05. manuell neu erarbeitet
    \item Git-Integration erst Ende Juni/Anfang Juli, UI entwickelte zeitweise auf eigenem Stand
    \item MVP-Meilenstein am 17.07. statt Ende Juli; Screen-Recordings statt Live-Demo, da Sophia und Nadine fehlten
    \item August ohne eigenen Meilenstein, Fokus auf Medien, Testgerät und Bericht
    \item Web-Build der App durch Isar eingeschränkt, separates Netlify-Vorschauprojekt
\end{itemize}

% Diagramm: Zeitstrahl als Abbildung, laut Vorgabe der Betreuerin vom 02.07. mit Zuordnung der Personen zu den Arbeitspaketen.

\subsection{Aufgabenverteilung}

Tabelle~\ref{tab:aufgabenverteilung} gibt einen Überblick darüber, welche
Aufgaben die einzelnen Gruppenmitglieder im Projekt übernommen und welche
Abschnitte des Berichts sie verfasst haben.

% Zuordnung aus Protokollen, Discord-Archiv und Commit-Daten.
% Spalte "Abschnitte im Bericht" wird beim Schreiben der Kapitel gefuellt,
% Quelle ist die Fu\ss zeile in report_bachelorprojekt_rehab.tex.

\begin{itemize}
    \item Rollen verschoben sich im Projektverlauf: Ilker wechselte zur Knowledge-Gruppe, Khaled, Ahmad und Nadine kamen zur UI-Gruppe
    \item Zwischenpräsentation und Endphase: Präsentation und Poster wurden auf die Gruppen aufgeteilt (04.09.)
    \item offene Punkte: Zuordnung für Parizoda Alimukhamedova und Tenzin Methok Taschi im Team klären
\end{itemize}

\renewcommand{\arraystretch}{1.3}
\begin{xltabular}{\textwidth}{@{} >{\raggedright\arraybackslash}p{4.4cm} >{\raggedright\arraybackslash}X >{\raggedright\arraybackslash}p{2.6cm} @{}}
    \caption{Aufgabenverteilung}
    \label{tab:aufgabenverteilung} \\
    \toprule
    \textbf{Name} & \textbf{Aufgaben im Projekt} & \textbf{Abschnitte im Bericht} \\
    \midrule
    \endfirsthead
    \toprule
    \textbf{Name} & \textbf{Aufgaben im Projekt} & \textbf{Abschnitte im Bericht} \\
    \midrule
    \endhead
    \midrule
    \multicolumn{3}{r}{\small Fortsetzung auf der nächsten Seite} \\
    \endfoot
    \bottomrule
    \endlastfoot
    Nadine Albers            & Logo und Wireframes (Figma), Übungsvideos, Mitarbeit am Bericht & \\
    Mika Schneider           & Related Work (Fokus Mensch), Ontologie und Übungsliste, Kapitel 8 & \\
    Louis Unruh              & Repository und Workflow, Knowledge Graph, Abschlusspräsentation & \\
    Parizoda Alimukhamedova  & \textit{im Team klären} & \\
    Khaled Almasri           & UI-Implementierung, App-Dokumentation & \\
    Amjad Faham              & UI-Implementierung, Netlify-Vorschau, Merge-Fix, UI-Vortrag & \\
    Daniel Fetzer            & API, Flutter- und Git-Setup, Plan- und Schmerz-Logik, Testgerät & \\
    Bjarne Kahrs             & API, Fragebogen, Reviews und Merge-Konflikte & \\
    Mohamad Kalaf            & Fragebogen-Team (30.04.), UI & \\
    Ilker Yasin Koc          & Digitalisierung und Auswertung der Befragung, MVP-Features, ab August Wissensbasis & \\
    Hoang Hiep Nguyen        & API, Discord-Organisation, Fragebogen, Profil- und Körper-Anbindung & \\
    Elif Özcelik             & Praxiskontakte, Therapeutenfragebogen, Auswertung, Wissensbasis, Kapitel 8 & \\
    Sophia Pichler           & Protokolle, Wireframes und Design, Interview, Praxiskontakte, MVP-Präsentation & \\
    Rabih Takouzli           & Related Work (Fokus Anwendung), UI, Bericht & \\
    Tenzin Methok Taschi     & \textit{im Team klären} & \\
    Ahmad Zahra              & UI-Implementierung (Körper, Videos), Sitzungsprotokoll am 04.09. & \\
    
\end{xltabular}
